% !PRINTED 382
\weberpage{382}
%% [running head, not body text -- commented out 2026-08-07] Page furniture (running
%% head, folio, signature) is excluded from this edition by standing policy.
%% 382 \hfill Fünfzehnter Abschnitt. \hfill .§ 104. % [sic] stray period before section mark as printed

Es gilt zunächst für jede beliebige Klasse $A$ und für beliebige
Stammteiler $\delta, \delta_1, \delta_2$:

\weq{7}{\chi(1, A) = 1,}
\weq{8}{\chi(\delta_1, A)\chi(\delta_2, A) = \chi(\delta_1\delta_2, A),}
\weq{9}{\chi(\delta, A_0) = 1,}

worin $\delta_1\delta_2$ das oben definierte symbolische Produkt ist und $A_0$
die Hauptklasse bedeutet, durch die die Zahl 1 darstellbar ist.

Sind für eine Klasse $A$ die Charaktere $\chi(\delta_1, A)$, $\chi(\delta_2, A)$ $\dots$
$\chi(\delta_\lambda, A)$ gegeben, so sind nach (4) und (8) \gesperrt{alle} $\chi(\delta, A)$ bestimmt.
Jeder dieser Charaktere kann aber $= +1$ oder $= -1$ sein,
und so ergeben sich $2^\lambda$ mögliche Bestimmungen über die Charaktere.
Diese sind aber nicht voneinander unabhängig.

Denn es folgt aus (5)
\[
Dy^2 \equiv (2ax + by)^2 \ (\mathrm{mod}\ 4m),
\]
und daraus folgt, wenn man $m$ relaliv prim zu $a$, folglich zu $y$
annimmt: % [sic] Gedruckt ist relaliv statt relativ.
\[
(Dy^2, m) = (\Delta, m) = 1.
\]
Daraus ergibt sich, daß für jede Klasse $A$

\weq{10}{\chi(\Delta, A) = +1.}

Nun läßt sich zu jedem Stammteiler $\delta$ ein bestimmter, von
$\delta$ verschiedener \gesperrt{komplementärer} Stammteiler $\delta'$ so bestimmen,
daß im Sinne der symbolischen Multiplikation

\weq{11}{\delta\delta' = \Delta}

ist. Das Komplement von $\delta'$ ist wieder $\delta$, und jeder Primteiler,
der in $\delta$ und $\delta'$ zugleich aufgeht, muß in $Q$ aufgehen.

Daraus folgt nach (10) für jede Klasse $A$:

\weq{12}{\chi(\delta, A) = \chi(\delta', A),}

und hiernach ist die Anzahl der möglichen Bestimmungen über
die $\chi(\delta, A)$ nur noch $2^{\lambda-1}$.

Man vereinigt in ein \gesperrt{Geschlecht (Genus)} alle Formenklassen $A$, in denen sämtliche Charaktere $\chi(\delta, A)$ übereinstimmen,
und gelangt zu dem Resultate:

5. \gesperrt{Die Anzahl der existierenden Geschlechter ist
höchstens gleich $2^{\lambda-1}$.}

Ob die Anzahl der existierenden Geschlechter wirklich so
groß ist, ist eine tiefere Frage, auf die wir später zurückkommen.
